\subsection*{先找回这一次购买的结论}

客户端缺少的第一项信息，是 K1 最后对应哪笔购买。只回答当前余额 7000，无法说明这 3000 是否由 K1 扣除；余额也可能因另一笔支出恰好变成 7000。只回答购买成功，又缺少后续查询和核对所需的订单、扣款编号。服务需要沿 S 找到已提交 P1，再返回 P1 的终态、O1、D1、商品、报价、金额以及这笔扣款完成时的余额。

这里把与一次已完成操作相联系、恢复时继续返回的业务内容称为稳定结果。稳定表示这些字段描述同一个已经作出的结论，包括成交当时确定的事实，以及业务拒绝时确定的原因。操作完成以后，订单仍会继续履约，账户仍有后续收支，商品目录也可能调价。这些变化应当有各自的查询位置，使一个字段不同时承担成交历史和当前状态两种含义。

沿用第四章的历史：P1 以 Q7 扣 3000 后，S2 使用另一份有效的 1000 分报价 Q9 完成购买，账户由 7000 变为 6000。迟到的 R2 仍按原 S 和输入取得 P1 的 \code{amount\_minor=3000} 与 \code{balance\_after\_minor=7000}；账户查询另用 \code{current\_balance\_minor=6000} 表示查询时现值。下面继续追问每个字段应从哪份历史资料读取。

\begin{figure}[htbp]
\centering
\begin{tikzpicture}[x=1cm,y=1cm]
\node[booknode] (a) at (0,0) {初始余额\\10000};
\node[booknode] (b) at (4,0) {P1 扣 3000\\当次余额 7000};
\node[booknode] (c) at (8.5,0) {P2 扣 1000\\当前余额 6000};
\draw[bookarrow] (a)--(b);
\draw[bookarrow] (b)--(c);
\node[bookresult] (d) at (3,-2) {恢复原 S：P1 / D1 / 7000};
\node[booknode] (e) at (8.5,-2) {查询账户：6000};
\draw[bookalert] (b)--(d);
\draw[bookarrow] (c)--(e);
\end{tikzpicture}
\caption{原扣款后的余额与账户当前余额分别来自两个时间依据。}
\label{fig:ch05-original-current}
\end{figure}


两个数值同时正确。D1 的 7000 说明 P1 在账户上的资金位置，6000 说明之后又有一笔 1000 分支出。如果恢复时重新读取 \code{accounts.balance\_minor} 填入 P1 的 \code{balance\_after\_minor}，同一 P1 就会先后返回 7000 与 6000。客户端保存的收据与后来恢复的收据发生差异，差异来源只是另一项合法操作。字段名中带有 after，也不足以提醒实现者选对来源；数据关系中应明确它来自 D1 的历史快照。

成交金额同样有时间依据。Q7 确定本次为 3000，后来目录涨至 3500，恢复继续返回 Q7 的成交快照。把当前商品表与旧购买临时连接，再直接选当前售价，会把旧结果改写为 3500。只保存 \code{product\_id} 而删除价格快照，服务未来也可能无法解释当时扣款依据。商品当前名称可供页面补充展示，收据上的商品描述若属于承诺内容，则应保存成交时版本或文本，不能默认由当前目录替代。

\subsection*{操作终态与订单现状分别变化}

订单 O1 后来可能从已创建进入已发货，甚至出现退款。购买创建操作 P1 的 \code{succeeded} 表示这次购买创建和本地扣款已经完成。随后履约会继续推进，订单可能发货，也可能按业务流程退款；这些后续动作都以原成交为依据。因此恢复 P1 时继续返回 \code{succeeded}，订单查询可以同时返回 O1 的已发货状态。若把 \code{orders.status} 直接写到 P1 的操作结果中，原本含义明确的购买成功会随履约变成另一个操作结论。

退款也有自己的事实。假设 D1 扣 3000 以后，经过合法退款流程产生返还记录 F1，金额为 3000。账户增加 3000，订单现状标记已退款，D1 仍表示原扣款。恢复 P1 仍返回原成功及 D1；查询订单则展示扣款与退款的关系。把原购买结果改成未扣款，会掩盖曾经发生的资金变化，也使已经收到原成功的客户端无法解释后来的记录。本章没有实现退款流程，这一例子只说明结果字段应引用哪个时间的事实。

可以让一个页面同时展示原成交与订单现状，不过响应应把它们分开，并说明观察时间。P1 的稳定内容保存在 \code{original\_result} 中。本次查询到的订单状态保存在 \code{current\_order} 中，其中还包含观察时间字段 \code{observed\_at}。服务用这个字段记录本次观察时点，帮助调用方判断资料新旧。它本身不证明两个查询来自同一数据库快照。如果页面要求余额与订单处于同一读取时点，应采用能够提供共同读取视图的事务安排；分别读取时，相同时间标签并不意味着这次观测是原子的。Read Committed 的两条普通查询即使放在同一事务中，也可能取得不同的语句快照，因此不能只凭开启事务就作出这项承诺。

\begin{figure}[htbp]
\centering
\begin{tikzpicture}[x=1cm,y=1cm]
\node[bookresult,text width=3cm] (p) at (0,0) {P1 原终态：成功\\D1 扣款 3000};
\node[booknode,text width=3cm] (o) at (5,0) {O1 已创建\\随后进入履约};
\node[booknode,text width=3cm] (f) at (5,-2) {发货或后续退款\\退款 F1 引用 D1};
\node[booknote,text width=4cm] (r) at (0,-2) {恢复 P1 继续返回原成功、原订单与原扣款编号};
\draw[bookarrow] (p)--(o);
\draw[bookarrow] (o)--(f);
\draw[bookalert] (p)--(r);
\end{tikzpicture}
\caption{订单履约与退款增加后续事实，原购买的已提交结论仍然保留。}
\label{fig:ch05-operation-order}
\end{figure}


本章为恢复入口选择更简单的表示：只返回稳定结果和本次通信信息。订单与账户现状由各自查询入口提供。这样恢复 P1 无须等待履约服务，也不会因当前商品资料暂时不可用而丢失本来已经保存的收据。若产品随后增加组合查询，可以沿上述字段边界扩展，不必修改 P1 的完成内容。

\subsection*{查询与重试的动作边界}

查询入口读取指定 S 的持久事实。它不会取得新身份、检查是否还能买一份、扣余额或补建订单。POST 重试则重新进入原操作协议：若已有终态，核对输入后返回；若原事务回滚且 S 尚未成立，它有机会通过唯一插入接着完成购买。因此相同 S 的查询与提交重试可能都找到 P1，但查不到时有不同后果。

设 R1 尚在路上，查询 Q1 先到数据库。Q1 没有找到记录，返回 404，表示这次读取没有找到目标记录；本接口还返回 \code{operation\_not\_found}，说明没有取得该身份的已提交终态。随后 R1 到达并提交 P1，账户变为 7000。这条历史中，404 与成功先后出现，两者并无矛盾。404 的证据范围只到查询当时和指定读取路径，不能取消路上的 R1。

若客户端此时把 K1 丢弃，换 K2 重发，R1 与 K2 就可以各买一份。正确继续方式是保留原身份和输入，再查询或使用同 S 重试提交。重试若遇到原事务未结束，唯一检查等待它的结局；原事务提交，返回原结果；原事务回滚，新尝试才可能取得 S。查询本身不做这项竞争，因此也不会通过查不到自动推动购买完成。

\begin{figure}[htbp]
\centering
\begin{tikzpicture}[x=1cm,y=1cm]
\node[booknote] at (0,.5) {客户端};
\node[booknote] at (8,.5) {服务与权威数据库};
\draw[draw=Muted] (0,0)--(0,-4.8);
\draw[draw=Muted] (8,0)--(8,-4.8);
\draw[bookarrow] (0,-.3)--(4,-.3)--(4,-3.3)--(8,-3.3);
\node[booknote,anchor=south] at (2,-.3) {R1：原 POST 已发出};
\node[booknote,anchor=west] at (4.2,-2.7) {网络延迟};
\draw[bookarrow] (0,-1)--node[above,font=\small,fill=white] {Q1 查询完整 S} (8,-1.3);
\draw[bookalert] (8,-1.8)--node[above,font=\small,fill=white] {404：此刻未找到} (0,-2.1);
\node[booknote,anchor=south] at (6,-3.3) {R1 此时才到达};
\node[bookresult] at (8,-4.3) {提交 P1 / O1 / D1};
\node[booknote,text width=5cm] at (1,-4.1) {Q1 没有撤销在途请求\\客户端继续保存原 S};
\end{tikzpicture}
\caption{一次查空之后，在途购买仍可能提交。}
\label{fig:ch05-query-before-post}
\end{figure}


查询还能用于客户端没有原响应、但保存了 S 的恢复。它无须重新提交完整购买体，因此不会通过查询判断一份新输入是否等于原输入；它只读取已经授权访问的指定操作。POST 携带输入，才需要按记录保存的规范化版本比较。如果应用想确认本地输入与服务记录相符，可以让查询返回允许展示的原输入摘要或字段，再作核对，但不把查询改为新的购买。

S 包含账户、动作、会话与键。查询路径里虽然只写 K1，其余部分仍必须来自已认证上下文、明确动作和 E1 许可。只按 K1 查找既可能读错账户，也可能把另一会话同名键当作本次结果。查询得到 P1 后，同样要检查当前调用者是否有权查看这笔购买。幂等键是定位依据，访问凭据才承担访问授权；知道一个难以猜测的键并不自动获得订单读取权。
\FloatBarrier
